\section*{Lời nói đầu}
\addcontentsline{toc}{section}{Lời nói đầu}

Các em học sinh thân mến,

\medskip
Tích phân là một trong những nội dung quan trọng của chương trình Giải tích lớp 12, nhưng phần lớn học sinh chỉ biết đến nó qua các công thức tính diện tích, thể tích mà chưa thực sự hiểu vì sao những công thức ấy lại đúng. Tài liệu này được biên soạn với mong muốn giúp các em nhìn thấy bản chất hình học ẩn sau tích phân, thay vì chỉ ghi nhớ và áp dụng một cách máy móc.

Điểm khác biệt của tài liệu so với nhiều sách tham khảo hiện nay là cách trình bày mỗi bài toán. Ở đây, các em sẽ không thấy công thức được đưa ra ngay từ đầu rồi theo sau là lời giải mẫu để học thuộc. Thay vào đó, mỗi nội dung đều bắt đầu bằng một tình huống hoặc câu hỏi để các em tự suy nghĩ, sau đó mới dần dần đi đến hướng giải quyết và lời giải đầy đủ.

Cụ thể, trong tài liệu các em sẽ gặp bốn loại hộp nội dung với mục đích riêng.

Hộp đầu tiên là hoạt động khám phá, thường mở đầu một bài học hoặc một nội dung mới bằng một tình huống gần gũi, có thể là một mảnh đất, một ổ bánh mì, hay một chiếc cổng chào. Khi gặp hộp này, các em đừng vội tìm công thức để áp dụng, mà hãy thử hình dung và suy đoán xem vấn đề đang được đặt ra là gì.

Hộp thứ hai là câu hỏi gợi mở, được trình bày với đường viền nét đứt. Đây là những câu hỏi nhỏ, dẫn dắt từng bước, giống như khi có một người ngồi cạnh gợi ý cho các em suy nghĩ. Các em nên tự trả lời những câu hỏi này trước khi đọc tiếp phần sau.

Hộp thứ ba là phân tích ý tưởng, nêu ra hướng giải quyết chung cho vấn đề nhưng chưa đi vào tính toán cụ thể. Đây là lúc thích hợp để các em tự đặt bút, phác thảo cách làm của riêng mình trước khi đối chiếu với lời giải.

Cuối cùng là lời giải, trình bày đầy đủ và chi tiết. Các em chỉ nên đọc phần này sau khi đã thử sức với ba phần trước đó.

Sở dĩ tài liệu được tổ chức theo trình tự như vậy là vì trong phòng thi, không có sẵn phần gợi ý hay phân tích ý tưởng cho các em dựa vào. Khả năng tự đặt câu hỏi trước một bài toán lạ, chẳng hạn như đồ thị cắt trục ở đâu, đại lượng nào có thể mang giá trị âm, mới là điều giúp các em xử lý tốt các bài toán mới trong đề thi, chứ không chỉ là việc nhớ công thức.

Bên cạnh bốn loại hộp trên, tài liệu vẫn giữ lại những hộp quen thuộc: hộp định lý nêu lại kiến thức được rút ra sau phần khám phá, hộp phương pháp giải trình bày quy trình chung cho một dạng toán, và hộp chú ý nhắc đến những sai lầm thường gặp mà các em nên tránh khi làm bài.

Hy vọng tài liệu sẽ là một người bạn đồng hành hữu ích trong quá trình tự học của các em.

\bigskip
\hfill \textit{Thầy Tâm dạy Toán}

\newpage
